% =====================================================================
%  CARNET D'ENTRAINEMENT — FICHE 6 — THÈME 2
%  Niveau MOYEN — Corrigé
% =====================================================================
\documentclass[11pt,a4paper]{article}
\usepackage[utf8]{inputenc}
\usepackage[T1]{fontenc}
\usepackage{lmodern}
\usepackage[french]{babel}
\usepackage{amsmath,amssymb,mathtools}
\usepackage{icomma}
\usepackage[version=4]{mhchem}
\usepackage{siunitx}
\sisetup{output-decimal-marker={,}}
\usepackage{xcolor}
\usepackage{colortbl}
\usepackage{tikz}
\usepackage[most]{tcolorbox}
\usepackage{enumitem}
\usepackage{array}
\usepackage{booktabs}
\usepackage{fancyhdr}
\usepackage[hidelinks]{hyperref}
\usetikzlibrary{arrows.meta,positioning,calc,shapes.geometric}
\usepackage[a4paper,margin=2.1cm,headheight=26pt,headsep=8pt]{geometry}

\definecolor{primary}{RGB}{146,95,160}
\definecolor{primarydark}{RGB}{92,52,104}
\definecolor{corcol}{RGB}{0,114,178}
\definecolor{astucecol}{RGB}{198,93,9}
\definecolor{atomO}{RGB}{220,55,45}
\definecolor{lonepair}{RGB}{120,94,200}
\definecolor{bondcol}{RGB}{60,60,70}

\tcbset{boxbase/.style={enhanced,breakable,boxrule=0.9pt,arc=2.5pt,
   left=3mm,right=3mm,top=2.2mm,bottom=2.2mm,
   fonttitle=\bfseries,coltitle=white,toptitle=0.6mm,bottomtitle=0.6mm}}
\newtcolorbox[auto counter]{cor}[1][]{boxbase,colback=corcol!5,colframe=corcol,
   title={\textsc{Corrigé}~\thetcbcounter\ifx\relax#1\relax\relax\else\textmd{ —\itshape\ #1}\fi}}
\newtcolorbox{astucebox}[1][]{boxbase,colback=astucecol!8,colframe=astucecol,
   title={\textsc{Astuce}\ifx\relax#1\relax\relax\else\textmd{ : \itshape #1}\fi}}

% -------- macros des quatre grandeurs --------
\newcommand{\nmax}{n_{\max}(e^-)}
\newcommand{\nv}{n_{v}(e^-)}
\newcommand{\nl}{n_{\ell}(e^-)}
\newcommand{\nnl}{n_{n\ell}(e^-)}

\pagestyle{fancy}\fancyhf{}
\fancyhead[L]{\small\textcolor{primarydark}{\textbf{Fiche 6 · Thème 2} — Méthode des 5 étapes et structure de Lewis}}
\fancyhead[R]{\small\textcolor{primarydark}{\textsc{Moyen · Corrigé}}}
\fancyfoot[C]{\small\thepage}
\renewcommand{\headrulewidth}{0.8pt}
\renewcommand{\headrule}{\hbox to\headwidth{\color{primary}\leaders\hrule height \headrulewidth\hfill}}
\setlength{\parindent}{0pt}\setlength{\parskip}{3pt}

\newcommand{\rep}{\textbf{\textcolor{corcol}{Réponse : }}}

\begin{document}
\begin{center}
{\Large\bfseries\color{primarydark}Thème 2 — Corrigé du niveau Moyen}
\end{center}
\medskip

\begin{cor}[l'ammoniac \ce{NH3}]
\textbf{Étape 1.} Atome central : \ce{N} ; $a=1$, $b=3$.\\
\textbf{Étape 2.} $\nmax = 8\times 1 + 2\times 3 = \mathbf{14}$.\\
\textbf{Étape 3.} Neutre ($q=0$) : $\nv = 5 + 3\times 1 = \mathbf{8}$.\\
\textbf{Étape 4.} $\nl = 14 - 8 = 6 \Rightarrow 6/2 = \mathbf{3}$ liaisons \ce{N-H} (simples).\\
\textbf{Étape 5.} $\nnl = 8 - 6 = 2 \Rightarrow 2/2 = \mathbf{1}$ doublet non liant, porté par
l'azote.\\
\rep \ce{N} central, $3$ liaisons \ce{N-H} et $1$ doublet libre sur \ce{N}. Octet de \ce{N} :
$3\times 2 + 2 = 8\,e^-$. \checkmark
\end{cor}

\begin{cor}[le difluorure de soufre \ce{SF2}]
\textbf{Étape 1.} Atome central : \ce{S} ; tous visent l'octet, $a=3$ (\ce{S} et $2\,\ce{F}$),
$b=0$.\\
\textbf{Étape 2.} $\nmax = 8\times 3 + 0 = \mathbf{24}$.\\
\textbf{Étape 3.} Neutre : $\nv = 6 + 2\times 7 = \mathbf{20}$.\\
\textbf{Étape 4.} $\nl = 24 - 20 = 4 \Rightarrow 4/2 = \mathbf{2}$ liaisons \ce{S-F} (simples).\\
\textbf{Étape 5.} $\nnl = 20 - 4 = 16 \Rightarrow 16/2 = \mathbf{8}$ doublets non liants.
Répartition : chaque \ce{F} complète son octet avec $3$ doublets libres ($2\times 3 = 6$), et il
reste $8 - 6 = 2$ doublets libres pour le soufre.\\
\rep $2$ liaisons \ce{S-F} ; $8$ doublets libres au total ($2$ sur \ce{S}, $3$ sur chaque \ce{F}).
Le soufre est du type \ce{AX2E2} (coudé).
\end{cor}

\begin{cor}[le dioxyde de carbone \ce{CO2}]
\textbf{Étape 1.} Atome central : \ce{C} ; $a=3$, $b=0$.\\
\textbf{Étape 2.} $\nmax = 8\times 3 = \mathbf{24}$.\\
\textbf{Étape 3.} Neutre : $\nv = 4 + 2\times 6 = \mathbf{16}$.\\
\textbf{Étape 4.} $\nl = 24 - 16 = 8 \Rightarrow 8/2 = \mathbf{4}$ liaisons. Le carbone n'est lié
qu'à \emph{deux} oxygènes : $4$ liaisons pour $2$ atomes $\Rightarrow$ \textbf{deux doubles
liaisons} \ce{C=O}.\\
\textbf{Étape 5.} $\nnl = 16 - 8 = 8 \Rightarrow 8/2 = \mathbf{4}$ doublets non liants, soit
$2$ sur chaque oxygène. Le carbone ne porte aucun doublet libre.\\
\rep structure \ce{O=C=O} : deux doubles liaisons, $4$ doublets libres ($2$ par \ce{O}).

\begin{center}
\begin{tikzpicture}[scale=1.0]
  \coordinate (C) at (0,0);
  \coordinate (Ol) at (-1.8,0);
  \coordinate (Or) at (1.8,0);
  \draw[bondcol,line width=1.3pt] ($(Ol)+(0.45,0.07)$)--($(C)+(-0.35,0.07)$);
  \draw[bondcol,line width=1.3pt] ($(Ol)+(0.45,-0.07)$)--($(C)+(-0.35,-0.07)$);
  \draw[bondcol,line width=1.3pt] ($(C)+(0.35,0.07)$)--($(Or)+(-0.45,0.07)$);
  \draw[bondcol,line width=1.3pt] ($(C)+(0.35,-0.07)$)--($(Or)+(-0.45,-0.07)$);
  \node[font=\large\bfseries] at (C) {C};
  \node[font=\large\bfseries,color=atomO] at (Ol) {O};
  \node[font=\large\bfseries,color=atomO] at (Or) {O};
  \fill[lonepair] ($(Ol)+(70:0.5)$) circle (2.2pt); \fill[lonepair] ($(Ol)+(110:0.5)$) circle (2.2pt);
  \fill[lonepair] ($(Ol)+(250:0.5)$) circle (2.2pt); \fill[lonepair] ($(Ol)+(290:0.5)$) circle (2.2pt);
  \fill[lonepair] ($(Or)+(70:0.5)$) circle (2.2pt); \fill[lonepair] ($(Or)+(110:0.5)$) circle (2.2pt);
  \fill[lonepair] ($(Or)+(250:0.5)$) circle (2.2pt); \fill[lonepair] ($(Or)+(290:0.5)$) circle (2.2pt);
\end{tikzpicture}
\end{center}
\end{cor}

\begin{cor}[la phosphine \ce{PH3}]
\textbf{Étape 1.} Atome central : \ce{P} ; $a=1$, $b=3$.\\
\textbf{Étape 2.} $\nmax = 8\times 1 + 2\times 3 = \mathbf{14}$.\\
\textbf{Étape 3.} Neutre : $\nv = 5 + 3\times 1 = \mathbf{8}$.\\
\textbf{Étape 4.} $\nl = 14 - 8 = 6 \Rightarrow 6/2 = \mathbf{3}$ liaisons \ce{P-H}.\\
\textbf{Étape 5.} $\nnl = 8 - 6 = 2 \Rightarrow \mathbf{1}$ doublet non liant sur le phosphore.\\
\rep $3$ liaisons \ce{P-H} et $1$ doublet libre sur \ce{P} : structure \emph{identique} à celle de
\ce{NH3} (\ce{P} et \ce{N} sont dans la même colonne, $N_v = 5$).
\end{cor}

\begin{cor}[le tétrachlorure de carbone \ce{CCl4}]
\textbf{Étape 1.} Atome central : \ce{C} ; $a=5$ (\ce{C} et $4\,\ce{Cl}$), $b=0$.\\
\textbf{Étape 2.} $\nmax = 8\times 5 = \mathbf{40}$.\\
\textbf{Étape 3.} Neutre : $\nv = 4 + 4\times 7 = \mathbf{32}$.\\
\textbf{Étape 4.} $\nl = 40 - 32 = 8 \Rightarrow 8/2 = \mathbf{4}$ liaisons \ce{C-Cl} (simples).\\
\textbf{Étape 5.} $\nnl = 32 - 8 = 24 \Rightarrow 24/2 = \mathbf{12}$ doublets non liants, soit
$3$ par chlore (le carbone n'en porte aucun).\\
\rep $4$ liaisons \ce{C-Cl} ; $12$ doublets libres ($3$ par \ce{Cl}). Type \ce{AX4} (tétraédrique).
\end{cor}

\begin{astucebox}[liaisons multiples : le déclic]
Si le nombre de liaisons ($\nl/2$) dépasse le nombre d'atomes \emph{voisins} de l'atome central,
c'est le signe qu'il y a des liaisons \textbf{multiples}. \ce{CO2} : $4$ liaisons pour $2$
voisins $\Rightarrow$ deux doubles liaisons. À l'inverse, pour \ce{CCl4} ($4$ liaisons, $4$
voisins) tout est simple.
\end{astucebox}

\end{document}
